% Generated draft: requires mathematical and source review.
\documentclass[11pt]{article}
\usepackage[margin=25mm]{geometry}
\usepackage{fontspec}
\setmainfont{Latin Modern Roman}
\usepackage{amsmath,amssymb,mathtools,unicode-math}
\setmathfont{Latin Modern Math}
\usepackage{enumitem,needspace,xurl,hyperref}
\setlength{\parindent}{0pt}
\setlength{\parskip}{5pt}
\setlength{\emergencystretch}{3em}
\newcommand{\N}{\mathbb{N}}
\newcommand{\Z}{\mathbb{Z}}
\newcommand{\Q}{\mathbb{Q}}
\newcommand{\R}{\mathbb{R}}
\newcommand{\C}{\mathbb{C}}
\newcommand{\F}{\mathbb{F}}
\newcommand{\A}{\mathbb{A}}
\newcommand{\PP}{\mathbb{P}}
\newcommand{\E}{\mathbb{E}}
\newcommand{\eps}{\varepsilon}
\newcommand{\dd}{\,\mathrm{d}}
\begin{document}
\section*{UnsolvedMath ID: 30005720}
\textbf{Source problem code:} catalogue entry
\subsection{Definitions and mathematical statement}
For a matroid \(M\), write \(E(M)\) for its ground set. A minor of \(M\) is a matroid obtained by deletion and contraction, considered up to isomorphism. The supplied record does not specify whether the matroids under consideration must be finite.

The term \emph{tensor product of matroids} is not defined in the supplied record. To display the logical structure without selecting an unsupported definition, let \(\mathcal M\) denote the intended class of matroids and let \(\mathsf T(M,N;P)\) be a placeholder for the assertion that \(P\) is a tensor product of \(M\) and \(N\) in the intended sense. Define, schematically,
\[
\mathcal C_{\mathsf T}
 =\bigl\{(M,N)\in\mathcal M\times\mathcal M:
       \exists P\in\mathcal M\;\mathsf T(M,N;P)\bigr\}.
\]
This notation supplies no tensor-product axioms and makes no assertion of existence or uniqueness.

\textbf{Problem, subject to the unresolved definitions.} Construct further examples of pairs in \(\mathcal C_{\mathsf T}\); give combinatorial descriptions of tensor products for particular classes of matroids; and characterize obstructions to membership in \(\mathcal C_{\mathsf T}\) by forbidden minors.

More explicitly, once the intended notion of tensor product and the relevant minor relation \(\preceq\) on pairs have been specified, a forbidden-minor characterization would seek a family \(\mathcal F\) of obstructing pairs such that, for every \((M,N)\in\mathcal M\times\mathcal M\),
\[
(M,N)\in\mathcal C_{\mathsf T}
\quad\Longleftrightarrow\quad
\nexists (A,B)\in\mathcal F\text{ with }(A,B)\preceq(M,N).
\]
This is a conditional formulation of the requested obstruction theory: the source does not specify that obstructions are pairs, define the relevant minor relation, establish the required closure under minors, or request a finite obstruction family. In particular, componentwise deletion and contraction must not be assumed to be the intended relation without checking the source. The missing tensor-product definition and obstruction convention prevent a fully specified mathematical statement.
\subsection{Short English statement}
Determine which pairs of matroids admit the intended tensor product, construct additional examples, describe products combinatorially for selected matroid classes, and develop a forbidden-minor characterization. The supplied record leaves the tensor-product definition, matroid domain, and obstruction convention unspecified, so this remains a schematic statement requiring source clarification.
\subsection{Sources}
\begin{itemize}
\item Local UnsolvedMath registry, permanent ID \(30005720\), title \emph{Tensor Products of Matroids}, source problem code \texttt{catalogue entry}; collection: \emph{Oberwolfach Reports Open Problems}. No individual proposer is supplied; the proposed year is recorded as 2023.
\item Supplied report citation: \emph{Geometric, Algebraic and Topological Combinatorics} (2024), \emph{Oberwolfach Reports}, pp.\ 3249--3318, DOI \texttt{10.4171/owr/2023/58}: \url{https://doi.org/10.4171/owr/2023/58}. The original statement includes a page marker 3311 and an unresolved reference to [1].
\item Anderson, \emph{Matroid Products in Tropical Geometry}: \url{https://arxiv.org/abs/2306.14771}. This is a supplied literature reference; its contents have not been independently checked here.
\item Berczi et al., \emph{Matroid Products via Submodular Coupling}: \url{https://arxiv.org/abs/2411.02197}. This is a supplied literature reference; its contents have not been independently checked here.
\item Supplied literature assessment, checked \textbf{2026-08-22}: status \emph{partially solved}. The supplied assessment states that realizable matroids and some minor-closed tensor-product classes are understood, while the requested combinatorial descriptions and forbidden-minor theory remain incomplete. This is supplied dated curation derived from the OWR material, not an independently verified literature assessment; time-sensitive status and cited claims require independent verification.
\item The registry labels the statement \texttt{corrected\_verified} and reports that it was reconstructed and checked against the official Oberwolfach report. These are supplied provenance claims, not verification performed for this draft.
\end{itemize}
\end{document}
